\chapter{第九单元}

\Needspace{12\baselineskip}

\section{第九单元素养评价（A）}

\Needspace{8\baselineskip}

\pdfbookmark[2]{原书第73页}{orig-73}

\subsection*{原书第73页}

% p073_s01
\BookQuestion{一、1}
学校体育社团开展1分钟单人跳绳达标练习，三位同学一分钟跳绳成绩分别为：小红112个、小丽105个、小芳119个。这三名同学平均每人一分钟跳绳（\hspace{1.8em}）个。\par
\begin{BookAnswer}
填\(112\)。以112个为基准，105个少7个、119个多7个，多、少抵消。\[\begin{aligned}&112+(7-7)\div3\\={}&112+0\div3\\={}&112+0\\={}&112\text{(个)}\end{aligned}\]移多补少后每人同为112个，总数不变。
\end{BookAnswer}

% p073_s02
\BookQuestion{一、2}
有两瓶果汁，甲瓶重500克，乙瓶重400克，从甲瓶倒（\hspace{1.8em}）克到乙瓶中，两瓶果汁一样重，此时两瓶果汁都是（\hspace{1.8em}）克。\par
\begin{BookAnswer}
依次填\(50\)、\(450\)。甲倒出的质量是两瓶质量差的一半：\[\begin{aligned}&(500-400)\div2\\={}&100\div2\\={}&50\text{(克)}\end{aligned}\]总质量不变，再求两瓶相等时每瓶的质量：\[\begin{aligned}&(500+400)\div2\\={}&900\div2\\={}&450\text{(克)}\end{aligned}\]
\end{BookAnswer}

% p073_s03
\BookQuestion{一、3}
气象台在某天的2时、8时、14时、20时测得的温度分别是\(8^\circ\mathrm C\)、\(15^\circ\mathrm C\)、\(24^\circ\mathrm C\)、\(17^\circ\mathrm C\)，由此估计这天的平均温度是（\hspace{1.8em}）\({}^\circ\mathrm C\)。\par
\begin{BookAnswer}
填\(16\)。用加法交换律、结合律配成两个32，再除以观测次数4。\[\begin{aligned}&(8+15+24+17)\div4\\={}&[(8+24)+(15+17)]\div4\\={}&[32+32]\div4\\={}&64\div4\\={}&16\text{(摄氏度)}\end{aligned}\]
\end{BookAnswer}

% p073_s04
\BookQuestion{一、4}
三个数的平均数是30，其中两个数分别是24和16，第三个数是（\hspace{1.8em}）。\par
\begin{BookAnswer}
填\(50\)。总和是3份30，先将24与16凑成40再减。\[\begin{aligned}&30\times3-24-16\\={}&30\times3-(24+16)\\={}&90-40\\={}&50\end{aligned}\]检验：\((24+16+50)\div3=30\)。
\end{BookAnswer}

% p073_s05
\BookQuestion{一、5}
投篮比赛，甲队有6人，共投中102个球，乙队有7人，共投中112个球。（\hspace{1.8em}）队的成绩好。\par
\begin{BookAnswer}
填甲。两队人数不同，不能只比较投中总个数，要用总个数除以各队人数。甲队\(102\div6=17\)(个)，乙队\(112\div7=16\)(个)；\allowbreak{}\(17>16\)，甲队平均每人投中更多，整体成绩更好。
\end{BookAnswer}

% p073_s06
\BookQuestion{一、6}
某手机旗舰店，5月上旬卖了98部手机，中旬卖了125部，下旬卖了149部，这个月平均每天卖（\hspace{1.8em}）部。\par
\begin{BookAnswer}
填\(12\)。五月有31天，“平均每天”应把全月销量除以31，不能除以三个旬。\[\begin{aligned}&(98+125+149)\div31\\={}&(223+149)\div31\\={}&372\div31\\={}&12\text{(部)}\end{aligned}\]
\end{BookAnswer}

% p073_s07
\BookQuestion{一、7}
甲、乙两个数的平均数是80，丙是50，这三个数的平均数是（\hspace{1.8em}）。\par
\begin{BookAnswer}
填\(70\)。\textbf{方法一：以80为基准移多补少}\par 以80为基准，甲、乙的总量正好是2份80；若丙也为80，三数平均仍为80。丙实际少30，应把这30平均从3份中扣去，所以先用30除以3得10，再从80减10，避免先求210。\[\begin{aligned}&80-(80-50)\div3\\={}&80-30\div3\\={}&80-10\\={}&70\end{aligned}\]原“先求总量再平均”的数量思路也成立，但不能把代表两份的80和代表一份的50直接平均。\par\textbf{方法二：先求总量再平均}\[\begin{aligned}&(80\times2+50)\div3\\={}&(160+50)\div3\\={}&210\div3\\={}&70\end{aligned}\]
\end{BookAnswer}

% p073_s08
\BookQuestion{一、8}
两根同样长的木板，每根长20厘米，粘连在一起，重复部分长4厘米。粘连后木板总长（\hspace{1.8em}）厘米。\par
\begin{BookAnswer}
填\(36\)。两块木板长度相加时，接头的\(4\)厘米被算了两次，实际只应计一次，所以\(20\times2-4=36\)(厘米)。另可先留整块，再加第二块不重叠的部分：\(20+(20-4)=36\)(厘米)。
\end{BookAnswer}

% p073_s09
\BookQuestion{一、9}
右图直条分别表示小明、小红和小华的身高，虚线（\hspace{1.8em}）的高度最接近他们的平均身高。\par
\BookFigure{figures/p073_s09-01.tex}
\begin{BookAnswer}
填③。把较高直条超过平均线的部分，移去填补较矮直条的不足，使三条一样高；移多补少前后总高度不变。由图形比较，③最接近平衡后的高度。②处上方多出的高度明显多于下方不足，④又高于所有直条中的两条；不能把某一人的身高直接当作平均身高。
\end{BookAnswer}

% p073_s10
\BookQuestion{一、10}
四个数的平均数是20，去掉一个数，剩下的三个数的平均数是18，去掉的这个数是（\hspace{1.8em}）。\par
\begin{BookAnswer}
填\(26\)。简便思路：留下3个数平均各比20少2，去掉的数应比20多3份2。\[\begin{aligned}&20+(20-18)\times3\\={}&20+2\times3\\={}&20+6\\={}&26\end{aligned}\]另可用总和之差，同时计算两个乘法：\[\begin{aligned}&20\times4-18\times3\\={}&80-54\\={}&26\end{aligned}\]
\end{BookAnswer}

\Needspace{8\baselineskip}

\pdfbookmark[2]{原书第74页}{orig-74}

\subsection*{原书第74页}

% p074_s01
\BookQuestion{二、1}
班级统计每位同学本学期课外阅读的本数，最多的读了42本，最少的读了28本，则全班平均每人读书数量可能是（\hspace{1.8em}）本。A.35；B.25；C.45；D.20。\par
\begin{BookAnswer}
选A。平均数反映整体水平，必须在最少的\(28\)本与最多的\(42\)本之间。只有\(35\)符合；\allowbreak{}\(25\)、\(20\)小于最小值，\allowbreak{}\(45\)大于最大值。本题可先用平均数的范围判断，无须知道每人的本数。
\end{BookAnswer}

% p074_s02
\BookQuestion{二、2}
四年级（1）班举行1分钟踢毽子比赛，男生的平均成绩是45下/分，女生的平均成绩是55下/分，下面说法正确的是（\hspace{1.8em}）。A.每个男生的成绩都比女生低；B.女生的整体情况好于男生；C.全班的平均次数一定是50下；D.男生平均成绩低，则男生中没有成绩超过女生平均水平的同学。\par
\begin{BookAnswer}
选B。女生平均\(55\)下，男生平均\(45\)下，所以女生整体成绩较好。平均数不代表每个人的具体成绩，不能推出“每个女生都高于每个男生”，也不能排除某个男生超过\(55\)下。求全班平均还须用全班总次数除以总人数，人数未知时，不能认定为\((45+55)\div2=50\)下。
\end{BookAnswer}

% p074_s03
\BookQuestion{二、3}
学校开设两个兴趣小组，四（1）班参加书法小组的有29人，参加机器人小组的有32人，两个兴趣小组都参加的有16人，四（1）班共有学生（\hspace{1.8em}）人。A.67；B.55；C.45；D.77。\par
\begin{BookAnswer}
两项人数相加时，两项都参加的\(16\)人被算了两次，要减去一次。因此至少参加一项的有\(29+32-16=45\)(人)。若全班每人至少参加其中一项，则选C。若有人两项都没参加，全班人数还须加上这部分人数。题面没有说明全员参与，不能无条件认定全班恰为\(45\)人。
\end{BookAnswer}

% p074_s04
\BookQuestion{二、4}
光明小学篮球队原有9名队员，平均体重是15千克。体重39千克的飞飞加入后，篮球队队员现在的平均体重和原来相比，（\hspace{1.8em}）。A.比原来轻；B.与原来一样；C.比原来重；D.无法确定。\par
\begin{BookAnswer}
按当前原始扫描的平均15千克，选C。因为\(39>15\)，加入后平均体重提高。\par 书后A与当前题面数据不符。若原平均体重应为45千克，则\(39<45\)，加入后平均体重降低，选A；实际原平均数需核实。
\end{BookAnswer}

% p074_s05
\BookQuestion{二、5}
学校购买了6台打印机，平均每台打印机的价格是500多元，那么这6台打印机的总价可能是（\hspace{1.8em}）元。A.2970；B.3300；C.3888；D.2000。\par
\begin{BookAnswer}
选B。“平均\(500\)多元”表示每台平均大于\(500\)元且小于\(600\)元，所以总价大于\(3000\)元且小于\(3600\)元，只有\(3300\)元符合。
\end{BookAnswer}

% p074_s06
\BookQuestion{二、6}
有三条不同长度的丝带，虚线处表示三条丝带的平均长度。想一想，第一条丝带的长度是（\hspace{1.8em}）cm。第二条20cm，第三条14cm，第三条的终点恰好在平均长度虚线上。A.8；B.14；C.17；D.6。\par
\BookFigure{figures/p074_s06-01.tex}
\begin{BookAnswer}
选A。第三条的终点恰落在平均线上，平均长度是14厘米。第二条多出的部分补给第一条，第一条长度为：\[\begin{aligned}&14-(20-14)\\={}&14-6\\={}&8\text{(厘米)}\end{aligned}\]
\end{BookAnswer}

% p074_s07
\BookQuestion{三、1（题图）}
下面是一条河各处的水深情况。（单位：厘米。）七处测量数据依次为70、120、170、240、190、130、60。警示牌：“水深危险，注意安全！”第（1）（2）问续原书75页。\par
\BookFigure{figures/p074_s07-01.tex}
\begin{BookAnswer}
七处水深分别为\(70,120,170,240,190,130,60\)厘米。原书75页三、1给出平均水深与过河风险的完整解答。
\end{BookAnswer}

\Needspace{8\baselineskip}

\pdfbookmark[2]{原书第75页}{orig-75}

\subsection*{原书第75页}

% p075_s01
\BookQuestion{三、1（1）}
原书74页河流水深图的七处水深依次为70、120、170、240、190、130、60厘米。（1）以这条河各处水深的平均数作为这条河的平均水深。请你算出这条河的平均水深。\par
\begin{BookAnswer}
\textbf{方法一：分组凑整求总和}\par 用加法交换律、结合律，把70与130、240与60配对。互相独立的分组在同一步计算。\[\begin{aligned}&(70+120+170+240+190+130+60)\div7\\={}&[(70+130)+(240+60)+(120+170+190)]\div7\\={}&[200+300+(290+190)]\div7\\={}&[500+480]\div7\\={}&980\div7\\={}&140\text{(厘米)}\end{aligned}\]答：这条河的平均水深是140厘米。\par\textbf{方法二：以100厘米为基准移多补少}\par 多出20、70、140、90、30，少30、40；利用加减运算的性质，把100与160、90减70分别括成一组。\[\begin{aligned}&100+[(70+30)+(20+140)+90-(30+40)]\div7\\={}&100+[100+160+90-70]\div7\\={}&100+[(100+160)+(90-70)]\div7\\={}&100+[260+20]\div7\\={}&100+280\div7\\={}&100+40\\={}&140\text{(厘米)}\end{aligned}\]答：这条河的平均水深是140厘米。
\end{BookAnswer}

% p075_s02
\BookQuestion{三、1（2）}
原书74页河流水深图的七处水深依次为70、120、170、240、190、130、60厘米。（2）小鸵鸟的身高是2米，它过河有危险吗？为什么？\par
\begin{BookAnswer}
每米是100厘米，2米有2份100厘米，按“每份数×份数”列式。先把身高化为厘米：\[\begin{aligned}2\text{米}&=(100\times2)\text{厘米}\\&=200\text{厘米}\end{aligned}\]答：有危险。最大水深\(240\)厘米\(>200\)厘米，深处超过身高。平均水深小于身高，不能保证每处水深都小于身高。
\end{BookAnswer}

% p075_s03
\BookQuestion{三、2}
星期六小雨去郊游，上午她从家到目的地走的是上坡路，用了12分钟，走了720米。下午按原路返回，只用了8分钟。他来回平均每分钟走了多少米？\par
\begin{BookAnswer}
\textbf{总路程括成整体，再除以总时间}\par 来回总路程按每程720米乘2程列式，并用括号表示整体。简算时利用\(20=2\times10\)和乘除运算的性质，先算\(2\div2\)，再除以10。\[\begin{aligned}&(720\times2)\div(12+8)\\={}&(720\times2)\div20\\={}&720\times(2\div2)\div10\\={}&720\times1\div10\\={}&720\div10\\={}&72\text{(米/分)}\end{aligned}\]答：来回平均每分钟走72米。
\end{BookAnswer}

% p075_s04
\BookQuestion{三、3（1）}
在单板滑雪U型池比赛中，表中是单板滑雪运动员小雪的评分。（1）7位裁判打分的平均分是多少？\par
\par\smallskip\Needspace{4\baselineskip}\begingroup\small\setlength{\tabcolsep}{3pt}
\noindent\begin{tabularx}{\linewidth}{|*{8}{>{\centering\arraybackslash}X|}}\hline
评委编号 & 1号 & 2号 & 3号 & 4号 & 5号 & 6号 & 7号\\\hline
成绩（分） & 80 & 87 & 83 & 85 & 88 & 97 & 82\\\hline
\end{tabularx}\par
\endgroup\smallskip
\begin{BookAnswer}
\textbf{方法一：以85分为基准移多补少}\par 87与83的偏差抵消，88与82的偏差抵消；只剩97比85多12、80比85少5。\[\begin{aligned}&85+(12-5)\div7\\={}&85+7\div7\\={}&85+1\\={}&86\text{(分)}\end{aligned}\]答：七位裁判评分的平均分是86分。\par \textbf{方法二：分组求总分再平均}\par 用加法交换律、结合律分组，同时计算独立的分组。\[\begin{aligned}&(80+87+83+85+88+97+82)\div7\\={}&[(87+83)+(88+82)+(80+85+97)]\div7\\={}&[170+170+(165+97)]\div7\\={}&[340+262]\div7\\={}&602\div7\\={}&86\text{(分)}\end{aligned}\]答：平均分是86分。
\end{BookAnswer}

% p075_s05
\BookQuestion{三、3（2）}
在单板滑雪U型池比赛中，表中是单板滑雪运动员小雪的评分。（2）如果去掉一个最高分和一个最低分后再计算，此时的平均分是多少？\par
\par\smallskip\Needspace{4\baselineskip}\begingroup\small\setlength{\tabcolsep}{3pt}
\noindent\begin{tabularx}{\linewidth}{|*{8}{>{\centering\arraybackslash}X|}}\hline
评委编号 & 1号 & 2号 & 3号 & 4号 & 5号 & 6号 & 7号\\\hline
成绩（分） & 80 & 87 & 83 & 85 & 88 & 97 & 82\\\hline
\end{tabularx}\par
\endgroup\smallskip
\begin{BookAnswer}
\textbf{方法一：从总分扣除最高、最低分}\par 去掉2个评分，总分与份数都要相应减少；分子先算602减97，分母同时算7减2。\[\begin{aligned}&(602-97-80)\div(7-2)\\={}&(505-80)\div5\\={}&425\div5\\={}&85\text{(分)}\end{aligned}\]答：去掉一个最高分和一个最低分后，平均分是85分。\par\textbf{方法二：移多补少}\par 从87分移2分给83分，从88分移3分给82分，原来的85分不动，五个评分都成85分。按每份85分、5份求平均，利用乘除运算性质先算5除以5。\[\begin{aligned}&85\times5\div5\\={}&85\times(5\div5)\\={}&85\times1\\={}&85\text{(分)}\end{aligned}\]答：平均分是85分。
\end{BookAnswer}

% p075_s06
\BookQuestion{三、3（3）}
在单板滑雪U型池比赛中，表中是单板滑雪运动员小雪的评分。（3）这两个结果哪个代表小雪的成绩最合理？为什么？\par
\par\smallskip\Needspace{4\baselineskip}\begingroup\small\setlength{\tabcolsep}{3pt}
\noindent\begin{tabularx}{\linewidth}{|*{8}{>{\centering\arraybackslash}X|}}\hline
评委编号 & 1号 & 2号 & 3号 & 4号 & 5号 & 6号 & 7号\\\hline
成绩（分） & 80 & 87 & 83 & 85 & 88 & 97 & 82\\\hline
\end{tabularx}\par
\endgroup\smallskip
\begin{BookAnswer}
答：第二个结果较合理。剔除一个最高分和一个最低分，能减小极端评分对代表性结果的影响。这里“较合理”针对评分规则而言，并不表示第一个平均数算错。
\end{BookAnswer}

\Needspace{8\baselineskip}

\pdfbookmark[2]{原书第76页}{orig-76}

\subsection*{原书第76页}

% p076_s01
\BookQuestion{三、4（1）}
学校对四年级的学生进行绘画能力抽查。三位老师的评分如下（单位：分）。（1）将上面的表格填写完整。\par
\par\smallskip\Needspace{4\baselineskip}\begingroup\small\setlength{\tabcolsep}{3pt}
\noindent\textbf{四（1）班学生绘画能力抽查表（单位：分）}\par\smallskip
\noindent\begin{tabularx}{\linewidth}{|*{7}{>{\centering\arraybackslash}X|}}\hline
 & 1号 & 2号 & 3号 & 4号 & 5号 & 6号\\\hline
王老师 & 90 & 88 & 95 & 65 & 76 & 86\\\hline
李老师 & 94 & 90 & 96 &  & 80 & 88\\\hline
刘老师 & 95 & 86 & 94 & 66 & 75 & 93\\\hline
平均得分 & 93 & 88 &  & 67 & 77 & 89\\\hline
\end{tabularx}\par
\endgroup\smallskip
\par\smallskip\Needspace{4\baselineskip}\begingroup\small\setlength{\tabcolsep}{3pt}
\noindent\textbf{四（2）班学生绘画能力抽查表（单位：分）}\par\smallskip
\noindent\begin{tabularx}{\linewidth}{|*{7}{>{\centering\arraybackslash}X|}}\hline
 & 1号 & 2号 & 3号 & 4号 & 5号 & 6号\\\hline
王老师 & 83 & 93 & 93 & 83 & 88 & 77\\\hline
李老师 & 88 & 92 & 95 & 86 &  & 80\\\hline
刘老师 & 93 & 91 &  & 83 & 92 & 74\\\hline
平均得分 & 88 & 92 & 95 &  & 90 & 77\\\hline
\end{tabularx}\par
\endgroup\smallskip
\begin{BookAnswer}
五个空依次为\(70,95,90,97,84\)。每名学生有3位老师评分，必须沿同一学生所在列取数。补未知评分可用移多补少，求平均可先配对凑整或取基准数；以下按表中五个空依次说明。\par 四（1）班4号：另外两次评分比平均67分别少2、1，李老师应比67多3；利用加法结合律先合并2与1。\[\begin{aligned}&67+(67-65)+(67-66)\\={}&67+2+1\\={}&67+(2+1)\\={}&67+3\\={}&70\text{(分)}\end{aligned}\]四（1）班3号：以95分为基准，96多1分，94少1分，移多补少恰好抵消。\[\begin{aligned}&95+[(96-95)-(95-94)]\div3\\={}&95+[1-1]\div3\\={}&95+0\div3\\={}&95+0\\={}&95\text{(分)}\end{aligned}\]四（2）班5号：88与92先凑180，原式已易算，保持。\[\begin{aligned}&90\times3-(88+92)\\={}&270-180\\={}&90\text{(分)}\end{aligned}\]四（2）班3号：一次95恰等于平均，另一次93少2，刘老师应比95多2。\[\begin{aligned}&95+(95-93)\\={}&95+2\\={}&97\text{(分)}\end{aligned}\]四（2）班4号：以83为基准，86多出的3分平均分给3份。\[\begin{aligned}&83+(86-83)\div3\\={}&83+3\div3\\={}&83+1\\={}&84\text{(分)}\end{aligned}\]填完后分别求该列三次评分的平均，检查与表中平均分相符。\par 答：四（1）班应补70分、95分；四（2）班应补90分、97分、84分。
\end{BookAnswer}

% p076_s02
\BookQuestion{三、4（2）}
每个学生的平均得分若92～100分为优，80～91分为良，68～79分为达标，68分以下为待达标，将抽查成绩整理在下面的统计表中。（单位：人。）\par
\par\smallskip\Needspace{4\baselineskip}\begingroup\small\setlength{\tabcolsep}{3pt}
\noindent\textbf{四（1）班学生绘画能力抽查表（单位：分）}\par\smallskip
\noindent\begin{tabularx}{\linewidth}{|*{7}{>{\centering\arraybackslash}X|}}\hline
 & 1号 & 2号 & 3号 & 4号 & 5号 & 6号\\\hline
王老师 & 90 & 88 & 95 & 65 & 76 & 86\\\hline
李老师 & 94 & 90 & 96 &  & 80 & 88\\\hline
刘老师 & 95 & 86 & 94 & 66 & 75 & 93\\\hline
平均得分 & 93 & 88 &  & 67 & 77 & 89\\\hline
\end{tabularx}\par
\endgroup\smallskip
\par\smallskip\Needspace{4\baselineskip}\begingroup\small\setlength{\tabcolsep}{3pt}
\noindent\textbf{四（2）班学生绘画能力抽查表（单位：分）}\par\smallskip
\noindent\begin{tabularx}{\linewidth}{|*{7}{>{\centering\arraybackslash}X|}}\hline
 & 1号 & 2号 & 3号 & 4号 & 5号 & 6号\\\hline
王老师 & 83 & 93 & 93 & 83 & 88 & 77\\\hline
李老师 & 88 & 92 & 95 & 86 &  & 80\\\hline
刘老师 & 93 & 91 &  & 83 & 92 & 74\\\hline
平均得分 & 88 & 92 & 95 &  & 90 & 77\\\hline
\end{tabularx}\par
\endgroup\smallskip
\par\smallskip\Needspace{4\baselineskip}\begingroup\small\setlength{\tabcolsep}{3pt}
\noindent\textbf{统计表（单位：人）}\par\smallskip
\noindent\begin{tabularx}{\linewidth}{|*{5}{>{\centering\arraybackslash}X|}}\hline
班级 & 优 & 良 & 达标 & 待达标\\\hline
四（1）班 &  &  &  & \\\hline
四（2）班 &  &  &  & \\\hline
\end{tabularx}\par
\endgroup\smallskip
\begin{BookAnswer}
统计对象是\(12\)名学生各自的平均分，不是\(36\)个老师评分。按\(92\sim100\)为优、\(80\sim91\)为良、\(68\sim79\)为达标、低于\(68\)为待达标逐个归类，每人只计一次。\par 四（\(1\)）班为\(93,88,95,67,77,89\)；四（\(2\)）班为\(88,92,95,84,90,77\)。\[\begin{array}{c|rrrr|r}\text{班级}&\text{优}&\text{良}&\text{达标}&\text{待达标}&\text{合计}\\\hline\text{四(1)班}&2&2&1&1&6\\\text{四(2)班}&2&3&1&0&6\\\text{总计}&4&5&2&1&12\end{array}\]行合计各为\(6\)，列合计之和也为\(12\)，不重不漏。边界\(92\)归优、\(68\)归达标；题面“\(68\)分以下”在此按不含\(68\)理解，避免与上一段重叠。\par 答：四（1）班优2人、良2人、达标1人、待达标1人；四（2）班优2人、良3人、达标1人、待达标0人。
\end{BookAnswer}

% p076_s03
\BookQuestion{三、4（3）}
结合统计表，分析一下两班学生绘画能力的情况。\par
\par\smallskip\Needspace{4\baselineskip}\begingroup\small\setlength{\tabcolsep}{3pt}
\noindent\textbf{统计表（单位：人）}\par\smallskip
\noindent\begin{tabularx}{\linewidth}{|*{5}{>{\centering\arraybackslash}X|}}\hline
班级 & 优 & 良 & 达标 & 待达标\\\hline
四（1）班 &  &  &  & \\\hline
四（2）班 &  &  &  & \\\hline
\end{tabularx}\par
\endgroup\smallskip
\begin{BookAnswer}
两班都抽查\(6\)人，可以比较同一等级的人数。两班优秀人数均为\(2\)人；四（\(2\)）班良好有\(3\)人，比四（\(1\)）班多\(1\)人，并且待达标为\(0\)人。结合多个等级分布，可判断本次抽查中四（\(2\)）班整体较好；仅看优秀人数相同，还不足以下这个结论。\par 答：在本次抽查中，四（2）班学生绘画能力的整体情况较好。
\end{BookAnswer}

\Needspace{12\baselineskip}

\section{第九单元素养评价（B）}

\Needspace{8\baselineskip}

\pdfbookmark[2]{原书第77页}{orig-77}

\subsection*{原书第77页}

% p077_s01
\BookQuestion{一、1}
蛋糕店某款蛋糕前6天的销量分别是16块、15块、20块、18块、14块、19块，如果第7天准备的数量与前6天的平均销量相同，那么应该做（\hspace{1.8em}）块。\par
\begin{BookAnswer}
填\(17\)。按加法交换律、结合律配成三对34块，每对2天。用乘除运算的性质先算3除以3，最后求34块分给2天的平均销量。\[\begin{aligned}&(16+15+20+18+14+19)\div6\\={}&[(16+18)+(15+19)+(20+14)]\div6\\={}&[34+34+34]\div6\\={}&[34\times3]\div(3\times2)\\={}&34\times(3\div3)\div2\\={}&34\times1\div2\\={}&34\div2\\={}&17\text{(块)}\end{aligned}\]
\end{BookAnswer}

% p077_s02
\BookQuestion{一、2}
某商场第一季度共销售电脑237台，第二季度平均每月销售电脑145台。上半年一共销售了（\hspace{1.8em}）台，上半年平均每月销售（\hspace{1.8em}）台。\par
\begin{BookAnswer}
依次填\(672\)、\(112\)。第二季度有3个月，按每月145台乘3个月求季度总量，再加第一季度的总销量。\[\begin{aligned}&237+145\times3\\={}&237+435\\={}&672\text{(台)}\end{aligned}\]上半年由两个季度组成，共3加3个月，再求月平均：\[\begin{aligned}&672\div(3+3)\\={}&672\div6\\={}&112\text{(台)}\end{aligned}\]
\end{BookAnswer}

% p077_s03
\BookQuestion{一、3}
5个连续自然数的平均数是26，这5个自然数中最小的数是（\hspace{1.8em}），最大的数是（\hspace{1.8em}）。\par
\begin{BookAnswer}
依次填\(24\)、\(28\)。连续五个自然数关于中间数对称：最小数与最大数的和、次小数与次大数的和，都等于中间数的\(2\)倍；移多补少后平均数就是中间数\(26\)。所以最小为\(26-2=24\)，最大为\(26+2=28\)。
\end{BookAnswer}

% p077_s04
\BookQuestion{一、4}
3个数的平均数是12，其中两个数的平均数是11，第3个数是（\hspace{1.8em}）。\par
\begin{BookAnswer}
填\(14\)。已知两个数比平均12合计少2，第三个数应比12多2。\[\begin{aligned}&12+(12-11)\times2\\={}&12+1\times2\\={}&12+2\\={}&14\end{aligned}\]另可用总和之差，同时计算两个乘法：\[\begin{aligned}&12\times3-11\times2\\={}&36-22\\={}&14\end{aligned}\]
\end{BookAnswer}

% p077_s05
\BookQuestion{一、5}
航模小组制作飞机模型，小亮试飞了56次，比小静多试飞8次，两人平均每人试飞（\hspace{1.8em}）次。\par
\begin{BookAnswer}
填\(52\)。把两人相差的8次分成两半，较多的一方减4就成为平均数。\[\begin{aligned}&56-8\div2\\={}&56-4\\={}&52\text{(次)}\end{aligned}\]另一方法先求总次数再平均：\[\begin{aligned}&[56+(56-8)]\div2\\={}&[56+48]\div2\\={}&104\div2\\={}&52\text{(次)}\end{aligned}\]
\end{BookAnswer}

% p077_s06
\BookQuestion{一、6（1）}
下面是学校两个〔首字不清〕球队队员的身高情况。（单位：厘米。）（1）A队队员的平均身高是（\hspace{1.8em}）厘米。\par
\par\smallskip\Needspace{4\baselineskip}\begingroup\small\setlength{\tabcolsep}{3pt}
\noindent\textbf{A队}\par\smallskip
\noindent\begin{tabularx}{\linewidth}{|*{6}{>{\centering\arraybackslash}X|}}\hline
队员 & 1号 & 2号 & 3号 & 4号 & 5号\\\hline
身高 & 156 & 154 & 151 & 159 & 160\\\hline
\end{tabularx}\par
\endgroup\smallskip
\par\smallskip\Needspace{4\baselineskip}\begingroup\small\setlength{\tabcolsep}{3pt}
\noindent\textbf{B队}\par\smallskip
\noindent\begin{tabularx}{\linewidth}{|*{6}{>{\centering\arraybackslash}X|}}\hline
队员 & 1号 & 2号 & 3号 & 4号 & 5号\\\hline
身高 & 150 & 159 & 157 & 160 & \\\hline
\end{tabularx}\par
\endgroup\smallskip
\begin{BookAnswer}
填\(156\)。\textbf{方法一：以155厘米为基准移多补少}\par 以155厘米为基准：156多1、154少1，151少4、159多4，两对多、少部分分别抵消。只剩160比155多5厘米，平均补到5人，每人多1厘米。\[\begin{aligned}&155+(160-155)\div5\\={}&155+5\div5\\={}&155+1\\={}&156\text{(厘米)}\end{aligned}\]\par\textbf{另法：分组求总身高再平均}\par 两对身高分别合为310厘米，与160厘米合计后除以5。\[\begin{aligned}&(156+154+151+159+160)\div5\\={}&[(156+154)+(151+159)+160]\div5\\={}&[310+310+160]\div5\\={}&[620+160]\div5\\={}&780\div5\\={}&156\text{(厘米)}\end{aligned}\]平均数反映整体水平，不表示每位队员都高156厘米。
\end{BookAnswer}

% p077_s07
\BookQuestion{一、6（2）}
下面是学校两个〔首字不清〕球队队员的身高情况。（单位：厘米。）（2）B队队员的平均身高是157厘米，那么5号队员的身高是（\hspace{1.8em}）厘米。\par
\par\smallskip\Needspace{4\baselineskip}\begingroup\small\setlength{\tabcolsep}{3pt}
\noindent\textbf{A队}\par\smallskip
\noindent\begin{tabularx}{\linewidth}{|*{6}{>{\centering\arraybackslash}X|}}\hline
队员 & 1号 & 2号 & 3号 & 4号 & 5号\\\hline
身高 & 156 & 154 & 151 & 159 & 160\\\hline
\end{tabularx}\par
\endgroup\smallskip
\par\smallskip\Needspace{4\baselineskip}\begingroup\small\setlength{\tabcolsep}{3pt}
\noindent\textbf{B队}\par\smallskip
\noindent\begin{tabularx}{\linewidth}{|*{6}{>{\centering\arraybackslash}X|}}\hline
队员 & 1号 & 2号 & 3号 & 4号 & 5号\\\hline
身高 & 150 & 159 & 157 & 160 & \\\hline
\end{tabularx}\par
\endgroup\smallskip
\begin{BookAnswer}
填\(159\)。由平均数还原总身高，已知四人先配150与160凑310，再减已知总和。\[\begin{aligned}&157\times5-(150+159+157+160)\\={}&157\times5-[(150+160)+(159+157)]\\={}&785-[310+316]\\={}&785-626\\={}&159\text{(厘米)}\end{aligned}\]检验：\((150+159+157+160+159)\div5=785\div5=157\)。书后150不符，不能取A队5号数据。
\end{BookAnswer}

% p077_s08
\BookQuestion{一、6（3）}
下面是学校两个〔首字不清〕球队队员的身高情况。（单位：厘米。）（3）如果A队的2号队员和B队的2号队员互换，这时（\hspace{1.8em}）队队员的平均身高高一些。\par
\par\smallskip\Needspace{4\baselineskip}\begingroup\small\setlength{\tabcolsep}{3pt}
\noindent\textbf{A队}\par\smallskip
\noindent\begin{tabularx}{\linewidth}{|*{6}{>{\centering\arraybackslash}X|}}\hline
队员 & 1号 & 2号 & 3号 & 4号 & 5号\\\hline
身高 & 156 & 154 & 151 & 159 & 160\\\hline
\end{tabularx}\par
\endgroup\smallskip
\par\smallskip\Needspace{4\baselineskip}\begingroup\small\setlength{\tabcolsep}{3pt}
\noindent\textbf{B队}\par\smallskip
\noindent\begin{tabularx}{\linewidth}{|*{6}{>{\centering\arraybackslash}X|}}\hline
队员 & 1号 & 2号 & 3号 & 4号 & 5号\\\hline
身高 & 150 & 159 & 157 & 160 & \\\hline
\end{tabularx}\par
\endgroup\smallskip
\begin{BookAnswer}
填A。交换A队154厘米与B队159厘米的队员后，两队仍各5人；A队总身高增加159减154，B队减少同样的量。先求人均增减，再求新平均。\par A队：\[\begin{aligned}&156+(159-154)\div5\\={}&156+5\div5\\={}&156+1\\={}&157\text{(厘米)}\end{aligned}\]B队：\[\begin{aligned}&157-(159-154)\div5\\={}&157-5\div5\\={}&157-1\\={}&156\text{(厘米)}\end{aligned}\]比较157大于156，交换后A队平均身高高一些。
\end{BookAnswer}

\Needspace{8\baselineskip}

\pdfbookmark[2]{原书第78页}{orig-78}

\subsection*{原书第78页}

% p078_s01
\BookQuestion{二、1}
周末，三个小朋友相约去玩“海盗船”，他们3人的平均身高是136cm。他们可能有（\hspace{1.8em}）人能玩“海盗船”。提示牌：“140cm以下儿童禁止乘坐！”A.0；B.0或1；C.0或1或2；D.1。\par
\BookFigure{figures/p078_s01-01.tex}
\begin{BookAnswer}
选C。三人总身高\(136\times3=408\)厘米。若3人都能乘坐，即使只按每人140厘米计，总身高也至少要\(140\times3=420\)厘米；\allowbreak{}\(408<420\)，所以不能3人全能乘坐。\par 身高为\((136,136,136)\)厘米时0人能坐；为\((146,131,131)\)厘米时1人能坐；为\((146,146,116)\)厘米时2人能坐。三组总身高都为408厘米，平均身高都为136厘米，故0、1、2人均可能。
\end{BookAnswer}

% p078_s02
\BookQuestion{二、2}
45名同学报名课外公益活动，但有10名同学因故没有参加任何活动。实际上，参加走进敬老院活动的有21人，参加垃圾分类宣传活动的有17人，两项公益活动都参加的有（\hspace{1.8em}）人。A.4；B.7；C.17；D.3。\par
\begin{BookAnswer}
选D。先扣除两项都未参加的\(10\)人，实际至少参加一项的有\(45-10=35\)人。把两项人数相加时，重复人数被多计一次，所以两项都参加的有\(21+17-35=3\)人。\par 也可分块核对：只参加第一项\(21-3=18\)人，只参加第二项\(17-3=14\)人；\allowbreak{}\(18+3+14+10=45\)，正好是全班人数。
\end{BookAnswer}

% p078_s03
\BookQuestion{二、3}
如果三个数的平均数是45，而且这三个数均大于40，那么最大的那个数可能是（\hspace{1.8em}）。A.42；B.70；C.49；D.60。\par
\begin{BookAnswer}
选C。三个数总和是3份45。其余两数都大于40，所以求最大数的上界：\[\begin{aligned}&45\times3-(40+40)\\={}&135-80\\={}&55\end{aligned}\]最大数应不小于平均数45，且严格小于55。49符合，例如49、43、43；42太小，60、70太大。
\end{BookAnswer}

% p078_s04
\BookQuestion{二、4}
甲、乙、丙三人谈论他们的平均年龄，甲说：“我比平均年龄多6岁。”乙说：“我比平均年龄少6岁。”下列说法正确的是（\hspace{1.8em}）。A.丙正好6岁；B.丙的年龄和平均年龄相等；C.乙不是最小的；D.无法确定丙的年龄。\par
\begin{BookAnswer}
B。甲比平均年龄多6岁，乙比平均年龄少6岁。把甲多出的6岁补给乙，两人都等于平均年龄，丙也必须等于平均年龄。因此选B。
\end{BookAnswer}

% p078_s05
\BookQuestion{二、5}
气温学上通常用一天中2时、8时、14时、20时4个时刻的气温的平均数作为一天的平均气温。右侧是某地区某四天中4个时刻的气温情况图，日平均气温大约是\(22^\circ\mathrm C\)的是（\hspace{1.8em}）。A.第一天；B.第二天；C.第三天；D.第四天。\par
\BookFigure{figures/p078_s05-01.tex}
\begin{BookAnswer}
选B。以22摄氏度的竖虚线为基准，第二天四个点高出与低于22摄氏度的部分大致相抵，移多补少后平均约为22摄氏度。第一天各点在线右侧；第三天各点在线左侧，第四天整体也偏左。原图没有精确刻度，不能凭空给各点标出具体温度。
\end{BookAnswer}

% p078_s06
\BookQuestion{二、6}
烟台市“非遗小传人”评选中，小兰秧歌和剪纸两项的平均分是94分，增加草编考核后，三项总分是281分，则小兰草编的得分是（\hspace{1.8em}）。A.93分；B.92分；C.94分；D.97分。\par
\begin{BookAnswer}
选A。94分是两项平均分，先按每项94分乘2项还原总分，再从三项总分281分中扣除。\[\begin{aligned}&281-94\times2\\={}&281-188\\={}&93\text{(分)}\end{aligned}\]
\end{BookAnswer}

\Needspace{8\baselineskip}

\pdfbookmark[2]{原书第79页}{orig-79}

\subsection*{原书第79页}

% p079_s01
\BookQuestion{三、1（1）}
光明小学为希望工程捐款。四（1）班有三个小组，捐款情况如下表。（1）四（1）班平均每组捐款多少元？\par
\par\smallskip\Needspace{4\baselineskip}\begingroup\small\setlength{\tabcolsep}{3pt}
\noindent\begin{tabularx}{\linewidth}{|*{4}{>{\centering\arraybackslash}X|}}\hline
组别 & 第一小组 & 第二小组 & 第三小组\\\hline
人数 & 9 & 11 & 10\\\hline
捐款总额/元 & 178 & 224 & 198\\\hline
\end{tabularx}\par
\endgroup\smallskip
\begin{BookAnswer}
\textbf{方法一：总额除以组数}\par 问的是“平均每组”，把总捐款平均分成3份。\[\begin{aligned}&(178+224+198)\div3\\={}&(402+198)\div3\\={}&600\div3\\={}&200\text{(元)}\end{aligned}\]答：平均每组捐款200元。\par \textbf{方法二：移多补少}\par 从第二组的224元中移22元给第一组，再移2元给第三组；总钱数不变。三组移补后的金额分别为：\[178+22=200\text{(元)}\]\[\begin{aligned}&224-22-2\\={}&202-2\\={}&200\text{(元)}\end{aligned}\]\[198+2=200\text{(元)}\]三组恰好一样多。答：平均每组捐款200元。
\end{BookAnswer}

% p079_s02
\BookQuestion{三、1（2）}
光明小学为希望工程捐款。四（1）班有三个小组，捐款情况如下表。（2）四（1）班平均每人捐款多少元？\par
\par\smallskip\Needspace{4\baselineskip}\begingroup\small\setlength{\tabcolsep}{3pt}
\noindent\begin{tabularx}{\linewidth}{|*{4}{>{\centering\arraybackslash}X|}}\hline
组别 & 第一小组 & 第二小组 & 第三小组\\\hline
人数 & 9 & 11 & 10\\\hline
捐款总额/元 & 178 & 224 & 198\\\hline
\end{tabularx}\par
\endgroup\smallskip
\begin{BookAnswer}
\textbf{求总金额和总人数}\par 问的是“平均每人”，用总捐款除以总人数。三组人数不同，不能把各组人均捐款直接平均。总金额与总人数的两条计算支路在每一步同时计算。\[\begin{aligned}&(178+224+198)\div(9+11+10)\\={}&(402+198)\div(20+10)\\={}&600\div30\\={}&20\text{(元)}\end{aligned}\]答：平均每人捐款20元。
\end{BookAnswer}

% p079_s03
\BookQuestion{三、2（1）}
一袋土豆共40个，任意取出5个分别称一称，结果如下表。（1）取出的这5个土豆，平均每个重多少克？\par
\par\smallskip\Needspace{4\baselineskip}\begingroup\small\setlength{\tabcolsep}{3pt}
\noindent\begin{tabularx}{\linewidth}{|*{6}{>{\centering\arraybackslash}X|}}\hline
编号 & 1 & 2 & 3 & 4 & 5\\\hline
质量/克 & 123 & 145 & 156 & 172 & 184\\\hline
\end{tabularx}\par
\endgroup\smallskip
\begin{BookAnswer}
\textbf{方法一：分组凑整后求平均}\[\begin{aligned}&(123+145+156+172+184)\div5\\={}&[(156+184)+(123+145+172)]\div5\\={}&[340+(268+172)]\div5\\={}&[340+440]\div5\\={}&780\div5\\={}&156\text{(克)}\end{aligned}\]答：抽出的土豆平均每个重156克。\par \textbf{方法二：以150克为基准移多补少}\par 多出6、22、34克，不足27、5克；先把6与34凑40；不足的27与5合为32。\[\begin{aligned}&150+[(6+34)+22-(27+5)]\div5\\={}&150+[40+22-32]\div5\\={}&150+[62-32]\div5\\={}&150+30\div5\\={}&150+6\\={}&156\text{(克)}\end{aligned}\]答：平均每个重156克。
\end{BookAnswer}

% p079_s04
\BookQuestion{三、2（2）}
一袋土豆共40个，任意取出5个分别称一称，结果如下表。（2）这袋土豆大约一共重多少千克？\par
\par\smallskip\Needspace{4\baselineskip}\begingroup\small\setlength{\tabcolsep}{3pt}
\noindent\begin{tabularx}{\linewidth}{|*{6}{>{\centering\arraybackslash}X|}}\hline
编号 & 1 & 2 & 3 & 4 & 5\\\hline
质量/克 & 123 & 145 & 156 & 172 & 184\\\hline
\end{tabularx}\par
\endgroup\smallskip
\begin{BookAnswer}
\textbf{用抽样平均质量估计整袋}\par 任意抽出的\(5\)个土豆平均每个\(156\)克，用这个平均质量代表整袋土豆的大致水平。按“每个质量\(\times\)个数\(=\)总质量”估计\(40\)个的质量。\[1\text{千克}=1000\text{克}\]\[\begin{aligned}&156\times40\\&=6240\text{(克)}\end{aligned}\]\[6240\text{克}=6.24\text{千克}\]答：这袋土豆大约重\(6.24\)千克。平均值不表示每个土豆恰好重\(156\)克，抽样估计也不等于称量整袋的精确结果。
\end{BookAnswer}

% p079_s05
\BookQuestion{三、3}
小刚前四次体育测试的平均成绩是89分，第五次体育测试的成绩是94分，他五次体育测试的平均成绩是多少？\par
\begin{BookAnswer}
\textbf{方法一：先求五次总分}\[\begin{aligned}&(89\times4+94)\div5\\&=(356+94)\div5\\&=450\div5\\&=90\text{(分)}\end{aligned}\]答：五次测试的平均成绩是\(90\)分。\textbf{方法二：以原平均分为基准}\par 若第五次也得\(89\)分，平均仍为\(89\)分；实际多\(5\)分，平均分给\(5\)次，每次多\(1\)分。\[\begin{aligned}&89+(94-89)\div5\\&=89+5\div5\\&=89+1\\&=90\text{(分)}\end{aligned}\]答：新的平均成绩是\(90\)分。
\end{BookAnswer}

\Needspace{8\baselineskip}

\pdfbookmark[2]{原书第80页}{orig-80}

\subsection*{原书第80页}

% p080_s01
\BookQuestion{三、4}
班级开展“绿色AI”活动，参加“AI环保创意赛”的有：李想、王晴、赵磊、周敏、刘浩、陈西；参加“低碳编程挑战赛”的有：王晴、赵磊、周敏、林芳、郭瑞。将下图补充完整，并求出参与“绿色AI”活动的一共有多少人。\par
\BookFigure{figures/p080_s01-01.tex}
\begin{BookAnswer}
只参加环保创意赛：李想、刘浩、陈西；两项都参加：王晴、赵磊、周敏；只参加编程挑战赛：林芳、郭瑞。填集合图时，三类姓名分别写在各自区域内部，每人只出现一次；两项都参加者必须写在重叠区。\textbf{方法一：相加后减去重复人数}\par 两张名单分别有\(6\)、\(5\)人，相加时中间的\(3\)人被数了两次，减去一次。\[\begin{aligned}&6+5-3\\&=11-3\\&=8\text{(人)}\end{aligned}\]答：一共有\(8\)人参加活动。\textbf{方法二：分成三块不重复的人数}\par 按“只参加环保、两项都参加、只参加编程”三部分计数，不重不漏。\[\begin{aligned}&(6-3)+3+(5-3)\\&=3+3+2\\&=6+2\\&=8\text{(人)}\end{aligned}\]答：一共有\(8\)人参加活动。
\end{BookAnswer}

% p080_s02
\BookQuestion{三、5（1）}
下面是四年级（2）班某次视力检查的情况。男生：4.9、5.1、4.4、5.0、4.8、5.2、4.1、5.0、4.7、4.9、5.1、4.6、5.2、4.2、4.8、5.3、5.0、5.1、4.5、5.2、5.0。女生：5.1、4.8、4.7、5.2、4.9、5.3、4.3、5.1、4.9、5.0、5.2、4.2、4.8、5.1、4.7、5.3、5.0、4.6、5.2、4.9、5.0、5.2、4.8。（1）将上面的数据进行整理，完成下面的统计表。\par
\par\smallskip\Needspace{4\baselineskip}\begingroup\small\setlength{\tabcolsep}{3pt}
\noindent\begin{tabularx}{\linewidth}{|*{5}{>{\centering\arraybackslash}X|}}\hline
视力 & 4.2及以下 & 4.3～4.6 & 4.7～4.9 & 5.0及以上\\\hline
男生/人 &  &  &  & \\\hline
女生/人 &  &  &  & \\\hline
总计/人 &  &  &  & \\\hline
\end{tabularx}\par
\endgroup\smallskip
\begin{BookAnswer}
先按性别分行，再依视力区间逐个划记，每人只计一次。\(4.2\)属于第一组，\allowbreak{}\(4.3\)、\(4.6\)属于第二组，\allowbreak{}\(4.7\)、\(4.9\)属于第三组，\allowbreak{}\(5.0\)属于第四组，不能把边界漏掉或算进两组。\[\begin{array}{c|rrrr|r}&\le4.2&4.3\sim4.6&4.7\sim4.9&\ge5.0&\text{合计}\\\hline\text{男生}&2&3&5&11&21\\\text{女生}&1&2&8&12&23\\\text{合计}&3&5&13&23&44\end{array}\]横向分别核对人数。男生：\[\begin{aligned}&2+3+5+11\\&=5+5+11\\&=10+11\\&=21\text{(人)}\end{aligned}\]女生：\[\begin{aligned}&1+2+8+12\\={}&(1+2)+(8+12)\\={}&3+20\\={}&23\text{(人)}\end{aligned}\]纵向按男女相加得到各段人数；两种方向得到的全班人数同为\(44\)，说明人数相符。\par 答：统计表按上表填写。
\end{BookAnswer}

% p080_s03
\BookQuestion{三、5（2）}
根据本页视力检查数据及统计表回答：（2）该班共有（\hspace{1.8em}）名学生；全班视力在（\hspace{1.8em}）范围内的人数最多，有（\hspace{1.8em}）人。\par
\par\smallskip\Needspace{4\baselineskip}\begingroup\small\setlength{\tabcolsep}{3pt}
\noindent\begin{tabularx}{\linewidth}{|*{5}{>{\centering\arraybackslash}X|}}\hline
视力 & 4.2及以下 & 4.3～4.6 & 4.7～4.9 & 5.0及以上\\\hline
男生/人 &  &  &  & \\\hline
女生/人 &  &  &  & \\\hline
总计/人 &  &  &  & \\\hline
\end{tabularx}\par
\endgroup\smallskip
\begin{BookAnswer}
依次填\(44\)、\(5.0\)及以上、\(23\)。\textbf{方法一：按性别合计}\[\begin{aligned}&21+23\\&=44\text{(人)}\end{aligned}\]答：全班共有\(44\)人。\textbf{方法二：按视力分段合计}\[\begin{aligned}&3+5+13+23\\&=8+13+23\\&=21+23\\&=44\text{(人)}\end{aligned}\]答：全班共有\(44\)人。\par 比较各分段总人数\(3,5,13,23\)，\allowbreak{}\(23\)最大，所以人数最多的范围是\(5.0\)及以上，共\(23\)人；不能只比较男生或女生那一行。
\end{BookAnswer}

% p080_s04
\BookQuestion{三、5（3）}
根据本页视力检查数据及统计表回答：（3）视力大于或等于5.0为正常，该班视力低于5.0的有（\hspace{1.8em}）人。根据本次体检结果，你想对班级同学说些什么？\par
\par\smallskip\Needspace{4\baselineskip}\begingroup\small\setlength{\tabcolsep}{3pt}
\noindent\begin{tabularx}{\linewidth}{|*{5}{>{\centering\arraybackslash}X|}}\hline
视力 & 4.2及以下 & 4.3～4.6 & 4.7～4.9 & 5.0及以上\\\hline
男生/人 &  &  &  & \\\hline
女生/人 &  &  &  & \\\hline
总计/人 &  &  &  & \\\hline
\end{tabularx}\par
\endgroup\smallskip
\begin{BookAnswer}
\textbf{方法一：总人数减正常人数}\[\begin{aligned}&44-23\\&=21\text{(人)}\end{aligned}\]答：视力低于\(5.0\)的有\(21\)人。\textbf{方法二：前三段直接相加}\[\begin{aligned}&3+5+13\\&=8+13\\&=21\text{(人)}\end{aligned}\]答：视力低于\(5.0\)的有\(21\)人。建议：保持正确读写姿势，注意用眼卫生，减少连续近距离用眼，保证户外活动时间，并按要求检查视力。
\end{BookAnswer}
